\section{Rare events and signed palettes}
\label{l:palettes}

An estimate for a typical partial permutation need not remain useful after
conditioning on a very rare mapping event.  The first result below explains
how a single retained event resolves this difficulty.  It converts a bound
on the accumulated conditional errors, valid at every retained setting,
into simultaneous bounds on many coset probabilities.  We then choose the
cosets by the representation being tested.  Both steps are statements about
finite probability spaces and permutation groups; no dynamics enter their
proofs.

Let $n=2h$ and partition $[n]=A\sqcup B$, with $|A|=|B|=h$.
A \emph{palette on $B$} is a partition $B=B_1\sqcup\cdots\sqcup B_u$
into nonempty color classes.  Its subgroup is
\[
 H=\prod_{i=1}^u S_{B_i}\le S_n,
\]
acting identically on $A$.  Put $l_i=|B_i|$, $p_i=l_i/h$, and
$\mathcal H(p)=-\sum_i p_i\log p_i$.  We use the analogous definitions
for palettes on $A$.  A left coset $xH$ prescribes each image in $A$
separately and prescribes the image set of every color class in $B$.

\subsection{An information bound valid under rare conditioning}

Here is the precise reveal experiment.  Let $(\Omega,P)$ be a finite
probability space and $g:\Omega\to S_n$ an endpoint map.  For each fixed
ordering $o=(c_1,\ldots,c_n)$ of the labels, let
$Y_1^o,\ldots,Y_n^o$ be observations on $\Omega$ such that $Y_j^o$
determines $g(c_j)$.  These observations may contain more information
than the endpoint.  Write $\mathcal F_r^o=\sigma(Y_1^o,\ldots,Y_r^o)$.
At a feasible history the previously observed endpoints leave exactly
$m=n-r$ available sites.  Suppose $\Delta_r^o\ge0$ is
$\mathcal F_r^o$-measurable and, for every subset $Q$ of these sites,
\begin{equation}
 P\{g(c_{r+1})\in Q\mid\mathcal F_r^o\}
 \le \frac{|Q|}{n-r}+n\Delta_r^o.
 \label{l:palette-reference-cap}
\end{equation}
All conditional statements concern histories of positive reference
probability.  Let $\mathbb E_{AB}$ denote averaging over independent
uniform orders within $A$ and $B$, with all labels of $A$ first;
define $\mathbb E_{BA}$ by reversing the two halves.

\begin{lemma}[Information from a retained reveal experiment]
\label{l:palette-information}
In the experiment just described, fix an integer $0\le v<h$ and an event
$\mathcal G\subseteq\Omega$ of reference probability $z>0$.
Assume that \eqref{l:palette-reference-cap} holds for both half-first
order conventions and every $0\le r<n-v$.  Suppose that every
$\omega\in\mathcal G$ satisfies
\begin{equation}
 \mathbb E_{AB}\sum_{r=0}^{n-v-1}\Delta_r^o(\omega)
 +\mathbb E_{BA}\sum_{r=0}^{n-v-1}\Delta_r^o(\omega)
 \le a_n.
 \label{l:palette-retained-cap}
\end{equation}
Let $\mu$ be the law of $g$ under $P(\,\cdot\mid\mathcal G)$.
For every palette on either half, its subgroup $H$, and every $x\in S_n$,
\begin{equation}
 \mu(xH)\le \frac{1}{[S_n:H]}
       \exp\{v\mathcal H(p)+n^2a_n-\log z\}.
 \label{l:palette-information-cap}
\end{equation}
In particular, if $a_n\le n^{-6}$ and $z=1-o(1)$, the last two terms
in the exponent are $o(1)$, uniformly over all palettes and cosets.
\end{lemma}

\begin{proof}
The event $\mathcal G$ and the retained endpoint law $\mu$ are fixed
before choosing the palette. Fix now a palette on $B$ and a coset $xH$,
and put $E=\mathcal G\cap\{g\in xH\}$. If $P(E)=0$ the conclusion
is immediate. Otherwise write $Q=P(\,\cdot\mid E)$. The three laws
are related by
\[
 P(E)=z\mu(xH),\qquad
 D(Q\Vert P)=-\log P(E).
\]
Choose the reveal order independently of the setting, using the $AB$
convention under both $P$ and $Q$. Since $E$ does not involve the order,
conditioning the setting on $E$ leaves this order uniformly distributed.
For a fixed order, data processing followed by the entropy chain rule
gives
\[
 -\log P(E)\ge
 \sum_{r=0}^{n-v-1}\mathbb E_Q
 D\bigl(\mathcal L_Q(Y_{r+1}^o\mid\mathcal F_r^o)
    \Vert\mathcal L_P(Y_{r+1}^o\mid\mathcal F_r^o)\bigr).
\]
For labels in $A$, membership in $xH$ prescribes a singleton target;
for a label in $B_i$ it prescribes the target set $xB_i$.
At a history that occurs under $Q$, let $a_r\ge1$ be the number of
unoccupied sites in this target.  The next endpoint lies in that target
with $Q$-probability one.  For any probability $R$ supported on an event
$F$ of positive reference probability, the identity
$D(R\Vert P')=D(R\Vert P'(\,\cdot\mid F))-\log P'(F)$
shows that its entropy is at least $-\log P'(F)$.
Apply this observation to the next observation and then use
\eqref{l:palette-reference-cap}.  Since $m=n-r\le n$,
\[
 \begin{split}
 -\log\left(\frac{a_r}{m}+n\Delta_r^o\right)
 &=\log(m/a_r)-\log\left(1+\frac{nm\Delta_r^o}{a_r}\right)\\
 &\ge\log(m/a_r)-n^2\Delta_r^o.
 \end{split}
\]
This lower bound is valid even when the displayed upper bound on the
reference probability exceeds one.  Averaging the chain rule over orders
therefore yields
\begin{equation}
 -\log P(E)\ge
 \mathbb E_{AB}\mathbb E_Q
       \sum_{r=0}^{n-v-1}\log((n-r)/a_r)-n^2a_n.
 \label{l:palette-rare-chain}
\end{equation}
The error bound uses \eqref{l:palette-retained-cap} pointwise on $E$.
It does not use an average error estimate under $P$ after the change
of law to $Q$.

It remains to count the information lost by omitting the last $v$ labels.
If all $n$ labels were revealed, the numerator counts would multiply
to $n!$, while the target counts would multiply to $\prod_i l_i!$.
The full logarithmic sum is consequently $\log[S_n:H]$, independently
of the compatible observations.  The last $v$ labels are in $B$.
For $v=0$ the missing contribution is zero.  For $v>0$, if their color
counts are $t_1,\ldots,t_u$, it is
\[
 \log\frac{v!}{\prod_i t_i!}
 \le v\mathcal H((t_i/v)_i).
\]
The inequality follows because the multinomial probability
of these counts under probabilities $t_i/v$ is at most one.
The last $v$ labels form a uniform subset of $B$ under the order average,
so $\mathbb E_{AB}(t_i/v)=p_i$.  Concavity of entropy gives expected
missing contribution at most $v\mathcal H(p)$.  Insert this bound in
\eqref{l:palette-rare-chain}, exponentiate, and use
$\mu(xH)=P(E)/z$.  The $BA$ order gives the same result for palettes on
$A$.  Every error term is independent of the chosen palette and coset,
which proves simultaneous uniformity.
\end{proof}

\subsection{Choosing a palette for a representation}

We next relate the entropy of a suitable palette to the dimension of a
representation.  We use the characteristic-zero branching rule
\cite[Theorem~9.2]{James1978} and the hook-length formula
\cite[Theorem~4.53]{etingof}.  The horizontal Pieri rule is the one-row
case of the induction rule in \cite[Section~16]{James1978}; its vertical
form follows by transposition and sign twist
\cite[Theorem~6.7]{James1978}.  Write $D_\gamma$ for the dimension of
the irreducible representation indexed by a partition $\gamma$.

\begin{lemma}[A dimension estimate in terms of the longest row or column]
There is an absolute $c_0>0$ such that, for all sufficiently large $b$,
\label{l:palette-degree}
if $\gamma\vdash b$, $L=\max(\gamma_1,\gamma'_1)$, and $k=b-L>0$,
then
\begin{equation}
 \log D_\gamma\ge
 \begin{cases}
 c_0 k\log(b/k),&k\le b/2,\\
 c_0 b,&k>b/2.
 \end{cases}
 \label{l:palette-degree-bound}
\end{equation}
In particular $\log D_\gamma\ge c_1k$ for an absolute $c_1>0$.
\end{lemma}

\begin{proof}
Transpose the diagram if necessary, so its first row has length $L$.
Apply the first-row hook bound \eqref{l:hook-rearrangement} with
row length $L$, $k$ lower boxes, and total size $b$.
The dimension of the lower diagram is at least one, so
\[
 \log D_\gamma\ge\log\binom bL-\frac{k(1+\log L)}{L}.
\]
If $k\le b/2$, the correction is $O(k\log b/b)$ and
$\binom bk\ge(b/k)^k$, proving the first case.  If
$b/10\le L<b/2$, the binomial coefficient has logarithm at least
$L\log(b/L)\ge(b/10)\log2$, and the correction is $O(\log b)$.
If $L<b/10$, every hook is at most $2L<b/5$, so
$D_\gamma\ge(b/e)^b/(b/5)^b=(5/e)^b$.
These bounds prove the second case and the stated consequence.
\end{proof}

\begin{lemma}[Signed palettes with controlled entropy]
\label{l:signed-palette}
There is an absolute $C_0<\infty$ such that, for every sufficiently
large $h$ and every $\beta\vdash h$, there are positive integers
$l_1,\ldots,l_u$ summing to $h$ and a product character
$\theta=\theta_1\otimes\cdots\otimes\theta_u$ of
$H=S_{l_1}\times\cdots\times S_{l_u}$, each $\theta_i$ either
trivial or sign, for which
\begin{equation}
 \operatorname{Hom}_H(\theta,V_\beta)\ne0,
 \qquad h\mathcal H((l_i/h)_i)\le C_0\log D_\beta.
 \label{l:signed-palette-bound}
\end{equation}
The same character occurrence holds for every placement of the classes
of these sizes on an $h$-element set.
\end{lemma}

\begin{proof}
Starting with $\beta$, remove a first row or a first column of maximal
length, and repeat on the remaining partition until it is empty.  The
successive lengths are the $l_i$.  Assign a trivial character to a row
removal and a sign character to a column removal.  If a first row is
removed, the remaining partition interlaces the original one: putting
back its size is a horizontal-strip addition in the Pieri rule.
Transposing gives a vertical-strip addition for a first column.
Iterated Pieri, transitivity of induction, and Frobenius reciprocity
give the character occurrence.  Conjugation supplies every placement.

Let $L=l_1$ and $k=h-L$.  Every hook in a strip of length $l$ at the
time of its removal is at most $2l$, because $l$ is the longer of the
current first row and first column.  Its hook product is at most
$(2l)^l\le(2e)^l l!$.  For the first strip use the sharper bound
$L!e^k$, obtained from $\log(1+x)\le x$ and the fact that the
sum of the numbers of boxes below the strip is $k$.
Remaining hooks are unchanged by each removal.  Thus, if
$M=h!/\prod_i l_i!$, the hook formula yields
\[
 \log M\le\log D_\beta+(1+\log(2e))k.
\]
To compare this multinomial count with entropy, set
$f(j)=j\log j-\log j!$, with $f(0)=0$.  Then $f(j)\ge0$ and
$0\le f(j)-f(j-1)\le1$.  Consequently
\[
 h\mathcal H((l_i/h)_i)-\log M
 =f(h)-\sum_i f(l_i)\le f(h)-f(L)\le k.
\]
If $k>0$, Lemma~\ref{l:palette-degree} absorbs this total $O(k)$
error into $C_0\log D_\beta$.  If $k=0$, the representation is
trivial or sign, and the one-color palette has entropy zero.
\end{proof}

\subsection{The transfer from both halves}

Fix once and for all
\begin{equation}
 0<\delta<\frac14,\qquad 2\delta C_0\le\frac1{10},
 \qquad v=\lfloor\delta n\rfloor.
 \label{l:palette-delta}
\end{equation}
The next theorem receives only coset caps.  In particular, it can be
applied to the output of Lemma~\ref{l:palette-information} without
referring again to its reveal experiment.

\begin{theorem}[Signed-palette transfer]
\label{l:palette-transfer}
Let $n$ tend to infinity through even integers.  For each $n$, partition
$[n]=A\sqcup B$ into equal halves and let $\mu$ be a probability on
$S_n$.  Suppose that $\varepsilon_n\to0$ and that, for every palette
on either half, every placement of its classes, its subgroup $H$, and
every $x\in S_n$,
\begin{equation}
 \mu(xH)\le [S_n:H]^{-1}
            \exp\{v\mathcal H(p)+\varepsilon_n\}.
 \label{l:palette-transfer-hyp}
\end{equation}
Here $\delta$ and $v$ satisfy \eqref{l:palette-delta}.  For all
sufficiently large $n$, uniformly over
$\lambda\notin\{(n),(1^n)\}$,
\begin{align}
 \|\widehat\mu(\lambda)\|_{\mathrm{HS}}^2
 &\le D_\lambda^{-1/3},\label{l:palette-fourier-hs}\\
 \|\widehat\mu(\lambda)\|_{\mathrm{op}}
 &\le D_\lambda^{-1/6}.
 \label{l:palette-fourier}
\end{align}
If in addition $|\widehat\mu(\mathrm{sgn})|=o(1)$, then
$\|\mu^{*4}-U_{S_n}\|_{\mathrm{TV}}=o(1)$, and consequently
$\|\mu^{*30}-U_{S_n}\|_{\mathrm{TV}}=o(1)$ as well.
Both conclusions also hold with $\mu$ replaced in every factor
by a probability $\mu_0$ satisfying
$\|\mu_0-\mu\|_{\mathrm{TV}}=o(1)$.
\end{theorem}

\begin{proof}
Put $J=\widehat\mu(\lambda)$ and $D=D_\lambda$, using a unitary model
$\rho_\lambda$.  We first bound $J$ on each half-isotypic component,
including its multiplicities.  Combining the halves will bound each
joint component; we then absorb the number of joint types into a small
power of $D$.  For a type $\beta$ of $S_B$, choose its signed
palette from Lemma~\ref{l:signed-palette}.  Equations
\eqref{l:palette-delta} and \eqref{l:palette-transfer-hyp} give,
for large $n$,
\begin{equation}
 \mu(xH)\le\frac{2D_\beta^{1/10}}{[S_n:H]}.
 \label{l:palette-selected-cap}
\end{equation}
This holds for every conjugate of the chosen subgroup within $S_B$.

Let $P_{H,\theta}$ be the orthogonal projection onto the
$\theta$-character space in the restriction to $H$.  Define the signed
mass $\theta_H$ to equal $\overline{\theta(h)}/|H|$ on $H$ and zero
elsewhere.  Right convolution by $\theta_H$ has Fourier
matrix $JP_{H,\theta}$.  On a coset $xH$ each resulting mass has absolute
value at most $\mu(xH)/|H|$.  Hence
\[
 |S_n|\sum_{x\in S_n}|(\mu*\theta_H)(x)|^2
 \le [S_n:H]\sum_{C\in S_n/H}\mu(C)^2
 \le 2D_\beta^{1/10}.
\]
Plancherel, retaining the $\lambda$ term, proves
\begin{equation}
 D\|JP_{H,\theta}\|_{\mathrm{HS}}^2\le2D_\beta^{1/10}.
 \label{l:palette-signed-hs}
\end{equation}

Let $\Pi_{B,\beta}$ denote the entire $\beta$-isotypic projection.
Average $P_{H,\theta}$ over conjugation by $S_B$.
On the $\beta$ carrier its rank is at least one, so Schur's lemma
makes this average at least $D_\beta^{-1}I$ on that carrier.
The group algebra acts as the identity on its multiplicity space.
The average is therefore at least
$D_\beta^{-1}\Pi_{B,\beta}$ in positive semidefinite order on
the full representation.  Tracing against $J^*J$ in
\eqref{l:palette-signed-hs} gives
\begin{equation}
 \|J\Pi_{B,\beta}\|_{\mathrm{HS}}^2
 \le \frac{2D_\beta^{11/10}}{D}.
 \label{l:palette-isotype-hs}
\end{equation}
The identical estimate holds for a type $\alpha$ of $S_A$.

Every constituent of the product of the two halves has a small carrier
on at least one side.  Indeed, the commuting projections
$\Pi_{A,\alpha}\Pi_{B,\beta}$ decompose the identity orthogonally,
and every nonzero such component has $D_\alpha D_\beta\le D$.
Use \eqref{l:palette-isotype-hs} for the smaller degree to obtain
\begin{equation}
 \|J\Pi_{A,\alpha}\Pi_{B,\beta}\|_{\mathrm{HS}}^2
 \le 2D^{-9/20}.
 \label{l:palette-pair-hs}
\end{equation}

We verify that the number of type pairs costs only a vanishing power
of $D$.  Let $L=\max(\lambda_1,\lambda'_1)$ and $k=n-L>0$.
Branching forces every half-type diagram to be contained in $\lambda$.
If the longest line is a column, transpose $\lambda$ and every contained
half-type diagram for this count.  Containment and dimensions are
preserved, and each resulting tail beyond the first row has at most
$k$ boxes.
The number of choices for either half-type is at most
$B_k=\sum_{j=0}^k p(j)$, where $p(j)$ is the number of partitions of
$j$.  The partition bound \eqref{l:partition-counts} gives
\[
 B_k\le(k+1)e^{3\sqrt k}=e^{O(\sqrt k)}.
\]
Lemma~\ref{l:palette-degree} implies
$\log B_k/\log D\to0$ uniformly: if $k\le\sqrt n$, its denominator
is at least $(c_0/2)k\log n$, and if $k>\sqrt n$, it is at least
$c_1k$.  Thus $B_k^2=D^{o(1)}$ uniformly over the nontrivial,
nonsign diagrams.  Orthogonality of the domain projections in
\eqref{l:palette-pair-hs} now yields
\[
 \|J\|_{\mathrm{op}}^2\le\|J\|_{\mathrm{HS}}^2
 \le 2B_k^2D^{-9/20}\le D^{-1/3}
\]
for sufficiently large $n$. This proves both
\eqref{l:palette-fourier-hs} and \eqref{l:palette-fourier}.

For four factors, Hilbert--Schmidt submultiplicativity, Plancherel
and Cauchy--Schwarz give
\[
 4\|\mu^{*4}-U_{S_n}\|_{\mathrm{TV}}^2
 \le |\widehat\mu(\mathrm{sgn})|^8
       +\sum_{\lambda\notin\{(n),(1^n)\}}D_\lambda^{1-4/3}
 =o(1),
\]
since the nonlinear sum is $Z_{1/3}(n)\to0$ by
Theorem~\ref{l:degree-all-powers}. Replacing the four factors by
$\mu_0$ costs at most $4\|\mu_0-\mu\|_{\mathrm{TV}}=o(1)$.
Convolution contracts total variation and preserves uniform measure,
so each four-factor conclusion implies its thirty-factor counterpart.

There is also a direct thirty-factor completion using only the
operator bound \eqref{l:palette-fourier}. It gives
\[
 4\|\mu^{*30}-U_{S_n}\|_{\mathrm{TV}}^2
 \le |\widehat\mu(\mathrm{sgn})|^{60}
       +\sum_{\lambda\notin\{(n),(1^n)\}}D_\lambda^{-8}.
\]
The sum tends to zero.  To see this directly, group the diagrams by
$k=n-\max(\lambda_1,\lambda'_1)$.  There are at most $2p(k)$
diagrams in each group.  Lemma~\ref{l:palette-degree} gives the
summable majorant $2e^{3\sqrt k-8c_1k}$, while each fixed $k>0$
has $D_\lambda\to\infty$ as $n\to\infty$.  Dominated convergence
applies.  Replacing the thirty factors one at a time changes total
variation by at most $30\|\mu_0-\mu\|_{\mathrm{TV}}$.
\end{proof}

The order of choices is now visible.  First choose the absolute
$C_0$ in Lemma~\ref{l:signed-palette}, then $\delta$ in
\eqref{l:palette-delta}.  A reveal construction supplies the reference
conditional bound \eqref{l:palette-reference-cap}, together with a
common event of mass $1-o(1)$ on which
\eqref{l:palette-retained-cap} holds with $a_n\le n^{-6}$ down to
$v$ unrevealed labels.  Lemma~\ref{l:palette-information}
supplies all the caps required by Theorem~\ref{l:palette-transfer}
with $\varepsilon_n=n^{-4}-\log z$.  The sign expectation remains a
separate scalar input.
